% Options for packages loaded elsewhere
\PassOptionsToPackage{unicode}{hyperref}
\PassOptionsToPackage{hyphens}{url}
\PassOptionsToPackage{dvipsnames,svgnames,x11names}{xcolor}
%
\documentclass[
  11pt,
]{article}
\usepackage{amsmath,amssymb}
\usepackage{setspace}
\usepackage{iftex}
\ifPDFTeX
  \usepackage[T1]{fontenc}
  \usepackage[utf8]{inputenc}
  \usepackage{textcomp} % provide euro and other symbols
\else % if luatex or xetex
  \usepackage{unicode-math} % this also loads fontspec
  \defaultfontfeatures{Scale=MatchLowercase}
  \defaultfontfeatures[\rmfamily]{Ligatures=TeX,Scale=1}
\fi
\usepackage{lmodern}
\ifPDFTeX\else
  % xetex/luatex font selection
  \setmainfont[]{Latin Modern Roman}
  \setmonofont[]{DejaVu Sans Mono}
\fi
% Use upquote if available, for straight quotes in verbatim environments
\IfFileExists{upquote.sty}{\usepackage{upquote}}{}
\IfFileExists{microtype.sty}{% use microtype if available
  \usepackage[]{microtype}
  \UseMicrotypeSet[protrusion]{basicmath} % disable protrusion for tt fonts
}{}
\makeatletter
\@ifundefined{KOMAClassName}{% if non-KOMA class
  \IfFileExists{parskip.sty}{%
    \usepackage{parskip}
  }{% else
    \setlength{\parindent}{0pt}
    \setlength{\parskip}{6pt plus 2pt minus 1pt}}
}{% if KOMA class
  \KOMAoptions{parskip=half}}
\makeatother
\usepackage{xcolor}
\usepackage[margin=2.2cm]{geometry}
\usepackage{longtable,booktabs,array}
\usepackage{calc} % for calculating minipage widths
% Correct order of tables after \paragraph or \subparagraph
\usepackage{etoolbox}
\makeatletter
\patchcmd\longtable{\par}{\if@noskipsec\mbox{}\fi\par}{}{}
\makeatother
% Allow footnotes in longtable head/foot
\IfFileExists{footnotehyper.sty}{\usepackage{footnotehyper}}{\usepackage{footnote}}
\makesavenoteenv{longtable}
\usepackage{graphicx}
\makeatletter
\def\maxwidth{\ifdim\Gin@nat@width>\linewidth\linewidth\else\Gin@nat@width\fi}
\def\maxheight{\ifdim\Gin@nat@height>\textheight\textheight\else\Gin@nat@height\fi}
\makeatother
% Scale images if necessary, so that they will not overflow the page
% margins by default, and it is still possible to overwrite the defaults
% using explicit options in \includegraphics[width, height, ...]{}
\setkeys{Gin}{width=\maxwidth,height=\maxheight,keepaspectratio}
% Set default figure placement to htbp
\makeatletter
\def\fps@figure{htbp}
\makeatother
\setlength{\emergencystretch}{3em} % prevent overfull lines
\providecommand{\tightlist}{%
  \setlength{\itemsep}{0pt}\setlength{\parskip}{0pt}}
\setcounter{secnumdepth}{5}
\newlength{\cslhangindent}
\setlength{\cslhangindent}{1.5em}
\newlength{\csllabelwidth}
\setlength{\csllabelwidth}{3em}
\newlength{\cslentryspacingunit} % times entry-spacing
\setlength{\cslentryspacingunit}{\parskip}
\newenvironment{CSLReferences}[2] % #1 hanging-ident, #2 entry spacing
 {% don't indent paragraphs
  \setlength{\parindent}{0pt}
  % turn on hanging indent if param 1 is 1
  \ifodd #1
  \let\oldpar\par
  \def\par{\hangindent=\cslhangindent\oldpar}
  \fi
  % set entry spacing
  \setlength{\parskip}{#2\cslentryspacingunit}
 }%
 {}
\usepackage{calc}
\newcommand{\CSLBlock}[1]{#1\hfill\break}
\newcommand{\CSLLeftMargin}[1]{\parbox[t]{\csllabelwidth}{#1}}
\newcommand{\CSLRightInline}[1]{\parbox[t]{\linewidth - \csllabelwidth}{#1}\break}
\newcommand{\CSLIndent}[1]{\hspace{\cslhangindent}#1}
\ifLuaTeX
\usepackage[bidi=basic]{babel}
\else
\usepackage[bidi=default]{babel}
\fi
\babelprovide[main,import]{american}
\ifPDFTeX
\else
\babelfont{rm}[]{Latin Modern Roman}
\fi
% get rid of language-specific shorthands (see #6817):
\let\LanguageShortHands\languageshorthands
\def\languageshorthands#1{}
\usepackage{float}
\usepackage{fvextra}
\usepackage{microtype}
\usepackage{etoolbox}
\usepackage{fancyhdr}
\AtBeginEnvironment{CSLReferences}{\small\setstretch{1.0}\interlinepenalty=10000}
\pagestyle{fancy}
\fancyhf{}
\fancyfoot[C]{\small Article 1 / V11 / Author-review copy \quad \thepage}
\renewcommand{\headrulewidth}{0pt}
\setlength{\emergencystretch}{3em}
\setlength{\tabcolsep}{4pt}
\floatplacement{figure}{H}
\ifLuaTeX
  \usepackage{selnolig}  % disable illegal ligatures
\fi
\IfFileExists{bookmark.sty}{\usepackage{bookmark}}{\usepackage{hyperref}}
\IfFileExists{xurl.sty}{\usepackage{xurl}}{} % add URL line breaks if available
\urlstyle{same}
\hypersetup{
  pdftitle={Reference-Set Decoupled Co-Adaptive Training for Physics-Informed Neural Networks},
  pdfauthor={Youssef EL MOUTEE},
  pdflang={en-US},
  colorlinks=true,
  linkcolor={Maroon},
  filecolor={Maroon},
  citecolor={Blue},
  urlcolor={teal},
  pdfcreator={LaTeX via pandoc}}

\title{Reference-Set Decoupled Co-Adaptive Training for Physics-Informed
Neural Networks}
\author{Youssef EL MOUTEE}
\date{6 September 2026 - V11 author-review copy}

\begin{document}
\maketitle

\setstretch{1.12}
\emph{Original article title retained. Author metadata and final human
approval are pending.}

\hypertarget{abstract}{%
\section*{Abstract}\label{abstract}}
\addcontentsline{toc}{section}{Abstract}

Adaptive collocation changes the training measure observed by global
loss controllers in physics-informed neural networks (PINNs). We
introduce and evaluate reference-set decoupled co-adaptive training
(ReCoA-PINN), which estimates global component weights on a fixed
reference set while optimizing the network on adaptive collocation
points. A controlled study compares this mechanism with coupled control,
proposal correction, kernel-density volume weighting and variational
residual attention on a common PyTorch backbone. Historical paired
experiments do not establish general stabilization or accuracy gains
from the fixed reference. We extend the evaluation with 200 V09
ablations and 560 V10 runs covering difficult observations, doubled
optimization budgets and common-loss comparisons. Independent Cole--Hopf
validation identifies a maximum discrepancy of 0.28578 in the historical
coarse Burgers reference. A Huber-versus-MSE confound is addressed
through new MSE-matched experiments. In the corrected comparison,
variational residual attention has the lowest median relative L2 error
in four observation regimes; 11 of 16 paired L2 contrasts pass Holm
correction, while no maximum-error contrast passes its separate
correction. These exploratory findings support configuration-dependent
conclusions. We distinguish the finite-population identity of fresh
importance-weighted samples from claims about the complete adaptive
training trajectory, and document the controls needed to assess
reference-set decoupling without assuming its superiority.

\textbf{Keywords:} physics-informed neural networks; adaptive
collocation; loss balancing; reference sets; importance weighting;
reproducibility.

\hypertarget{introduction}{%
\section{Introduction}\label{introduction}}

Physics-informed neural networks (PINNs) approximate a physical field
while penalizing the residual of its governing differential equation.
Their training objective generally combines the physics residual,
initial and boundary conditions, and, when available, observations.
Because these components can have different scales and convergence
rates, adaptive global loss balancing is frequently used. Residual-based
adaptive sampling is also used to allocate more collocation points to
poorly resolved regions.

Each mechanism is reasonable in isolation, but their joint use raises a
subtle question. If the global loss controller estimates the importance
of the physics term on the same nonuniform points selected by the
residual sampler, it no longer observes the original uniform-domain
objective. The controller and the sampler may therefore reinforce one
another. This interaction is distinct from the well-known imbalance
between loss components: it concerns the measure under which a component
is evaluated.

Previous work has already established adaptive loss balancing,
residual-based sampling, importance sampling, volume correction and
variational residual attention. Therefore, our aim is not to claim the
invention of these individual mechanisms. We conduct a controlled study
of competing training measures: retaining the coupled controller,
estimating global weights on a fixed reference set, correcting a known
discrete proposal, estimating local volume by KDE, or deliberately
tilting the objective through vRBA.

The study asks five questions:

\begin{enumerate}
\def\labelenumi{\arabic{enumi}.}
\tightlist
\item
  Does residual sampling measurably alter the component losses and
  gradients observed by a global controller?
\item
  Does a fixed, independent reference set improve training stability or
  solution accuracy?
\item
  Does exact proposal correction improve a matched uncorrected adaptive
  sampler under sparse and noisy data?
\item
  Are the conclusions robust to PDE family and to an established
  alternative, KDE volume weighting?
\item
  Does a deliberately tilted variational objective improve global and
  local errors consistently on the same backbone?
\end{enumerate}

Our principal finding is conditional. The feedback is measurable and
exact correction can be beneficial. vRBA performs favorably in
historical relative-L2 comparisons, but a common-loss audit qualifies
the noisy-data claim; maximum error and cost do not show universal
dominance. This mixture of positive, null and negative results
identifies the controls required before an adaptive PINN method can
credibly claim robustness.

\hypertarget{related-work-and-positioning}{%
\section{Related work and
positioning}\label{related-work-and-positioning}}

Raissi, Perdikaris and Karniadakis established the modern PINN
formulation for forward and inverse differential-equation problems
(Raissi, Perdikaris, and Karniadakis 2019). Gradient-flow pathologies
and adaptive balancing were subsequently analyzed by Wang, Teng and
Perdikaris and by several self-adaptive weighting methods (Wang, Teng,
and Perdikaris 2021). Wu et al.~provided a broad comparison of
nonadaptive and residual-based adaptive collocation (Wu et al. 2023).
Nabian, Gladstone and Meidani introduced importance sampling for more
efficient PINN training (Nabian, Gladstone, and Meidani 2021).

The closest antecedent to the measure question is VW-PINNs by Song et
al.~They show that nonuniform collocation can make the standard
empirical PDE loss ill-conditioned and propose weighting residuals by
local volumes estimated via kernel density estimation (Song et al.
2025). Recent work has gone further: joint weighting and sampling
methods explicitly combine the two mechanisms, while the vRBA framework
of Toscano et al.~interprets residual adaptivity through convex
potentials and variational objectives (Toscano et al. 2026).

These studies prevent two overly broad claims: combining loss weighting
and adaptive sampling is not new, and correcting the effect of a
nonuniform collocation measure is not new. The potentially useful
contribution here is a paired controlled design that contrasts a fixed
controller reference, a known- proposal correction, KDE volume
correction, a variational tilted objective and matched controls under
exactly the same PyTorch backbone.

Chen, Howard and Stinis combine residual-decay weighting with adaptive
sampling (Chen, Howard, and Stinis 2025). Singh et al.~combine smoothed
gradient balancing and residual-based collocation on Burgers and
Allen--Cahn (Singh et al. 2026). These are direct antecedents to a broad
co-adaptation claim. Our M5 comparator is not a faithful reproduction of
either full published protocol. The present contribution is the explicit
test of a fixed controller reference together with matched mechanism
controls and documented limitations.

Modern point-placement alternatives include PINNACLE, which selects
collocation and experimental points jointly (Lau et al. 2024),
QR-DEIM-based collocation selection (Celaya, Fuentes, and Riviere 2025),
and gradient-based point movement in PACMANN (Visser, Heinlein, and
Giovanardi 2025). They are discussed but not evaluated here.
Consequently this study does not establish state-of-the-art superiority
over contemporary placement methods.

\hypertarget{methods}{%
\section{Methods}\label{methods}}

\hypertarget{base-objective}{%
\subsection{Base objective}\label{base-objective}}

For model parameters \(\theta\), the total objective is

\begin{equation}
\mathcal L(\theta)=
\sum_{j\in\{f,i,b,d\}}\lambda_j\mathcal L_j(\theta),
\end{equation}

where \(f\), \(i\), \(b\), and \(d\) denote the physics,
initial-condition, boundary-condition and observation terms. The
standard physics term on collocation points \(z_i\) is

\begin{equation}
\mathcal L_f=\frac1N\sum_{i=1}^{N}r_\theta(z_i)^2.
\end{equation}

\hypertarget{controlled-variants}{%
\subsection{Controlled variants}\label{controlled-variants}}

M0 uses uniform collocation and fixed component weights. M1 adds
adaptive component weights. M3 adds residual sampling but retains fixed
component weights. M5 combines both mechanisms and computes the
controller update on the adaptive training set. ReCoA-PINN (M6) uses the
same adaptive-sampling protocol as M5 but computes the controller update
on a fixed independent audit set of equal cardinality.

M7U samples with replacement from a finite candidate pool according to
the residual proposal \(q\), but leaves the physics loss uncorrected. M7
uses the same sampling mechanism and random-number streams, and attaches
the Hansen--Hurwitz weight

\begin{equation}
w_i=\frac{1}{N_c q_i},
\qquad
\widehat{\mathcal L}_{f,\mathrm{M7}}
=\frac1N\sum_iw_i r_\theta(z_i)^2,
\end{equation}

where \(N_c\) is the candidate-pool size. Retained points keep the
proposal weight assigned when they were sampled. A uniform proposal
floor bounds the largest possible importance weight.

For VW and VW-CA, a Gaussian KDE estimates the density of the realized
collocation set after mapping all coordinates to the unit hypercube.
With \(V_i\propto1/\widehat p_h(z_i)\), the physics loss is

\begin{equation}
\mathcal L_{f,\mathrm{VW}}
=\frac{\sum_iV_i^2r_\theta(z_i)^2}{\sum_iV_i^2}.
\end{equation}

VW uses fixed global component weights. VW-CA uses the same adaptive
global controller as M5 and M7. The default KDE bandwidth follows
Scott's rule; separate runs use 0.5, 1 and 2 times this bandwidth.

For vRBA, fixed point sets receive persistent local multipliers. With
the exponential potential, the numerically stable normalized attention
target is

\begin{equation}
\widehat q_i^k=
\exp\!\left(\frac{|r_i|-\max_j|r_j|}{\epsilon_k}\right),
\qquad
\epsilon_k=\frac{c\max_j|r_j|}{\log(k+2)}.
\end{equation}

The local update and component loss are

\begin{equation}
\lambda_i^{k+1}=\gamma_k\lambda_i^k+
\eta\left(\phi\widehat q_i^k+1-\phi\right),
\qquad
\mathcal L_j^{\mathrm{vRBA}}=
\frac1{N_j}\sum_i(\lambda_{j,i}r_{j,i})^2.
\end{equation}

The physics component keeps global weight one; other component weights
track ratios of exponentially smoothed gradient norms. We use the
published first-order defaults (\texttt{eta=0.01}, \texttt{phi=0.8},
local cap schedule, gradient EMA \texttt{0.99}, global-weight EMA
\texttt{0.99975}). This is a controlled PyTorch common-backbone
implementation of the vRBA mechanism, not an exact reproduction of the
authors' JAX architecture, hard constraints, optimizer sequence or
training horizon.

\hypertarget{fixed-reference-controller-and-interpretation}{%
\subsubsection{Fixed-reference controller and
interpretation}\label{fixed-reference-controller-and-interpretation}}

Let \(R\) be the fixed reference set and \(S_k\) the adaptive
collocation set. ReCoA-PINN computes controller losses and gradient
norms on \(R\) while updating the network using \(S_k\). M5 uses \(S_k\)
for both. The physics and known-condition reference points are drawn
independently at initialization; observed sensors and their targets are
shared. The reference set is used repeatedly to control training and is
therefore not an independent test set for the trained model. Adaptive
trajectories can subsequently diverge even with identical random seeds.

For component \(j\), controller updates implement

\begin{equation}
g_{j,k}=\|\nabla_\theta L_j\|_2,\quad
\bar g_{j,k}=\beta\bar g_{j,k-1}+(1-\beta)g_{j,k},\quad
a_{j,k}=\frac{\operatorname{mean}_\ell\bar g_{\ell,k}}{\bar g_{j,k}}
\left(\frac{L_{j,k}}{L_{j,0}}\right)^\alpha.
\end{equation}

The implementation normalizes \(a_{j,k}\) to unit mean, clips it between
\(w_{\min}\) and \(w_{\max}\), renormalizes the clipped targets to unit
mean and then applies an EMA to the component weights. Here
\(\beta=0.95\), \(\alpha=0.5\), \(w_{\min}=0.1\) and \(w_{\max}=10\).
Clipping occurs before the second normalization: these settings are not
strict final-weight bounds. Positive numerical floors prevent division
by zero. This paper evaluates the implemented rule; it does not claim a
convergence theorem for the controller.

At each iteration, the algorithm (i) refreshes the scheduled fraction of
collocation points if due, (ii) updates global component weights from
the method-specific reference if due, (iii) constructs the weighted
training loss, and (iv) takes one Adam step with gradient clipping.
Evaluation points are never used for these updates. For vRBA, local
attention and its global controller are updated on fixed point sets; its
ablations are not ablations of adaptive point placement or of the M6
reference mechanism.

\hypertarget{scope-of-exact-proposal-correction}{%
\subsubsection{Scope of exact proposal
correction}\label{scope-of-exact-proposal-correction}}

For a frozen model \(\theta\) and candidate population \(C\), a fresh
with-replacement draw \(I\sim q\) satisfies

\begin{equation}
\mathbb E_{I\sim q}\left[\frac{r_\theta(z_I)^2}{N_c q_I}\;\middle|\;\theta,C,q\right]
=\frac{1}{N_c}\sum_{i=1}^{N_c}r_\theta(z_i)^2.
\end{equation}

This is a finite-population quadrature identity, not a guarantee for the
full adaptive optimization path. Retained points carry weights from
earlier candidate populations, and later model parameters depend on
earlier samples. The identity does not prove conditional unbiasedness of
every later gradient or elimination of continuous-domain discretization
error. M7 and M7U share the sampling algorithm and seeded random
streams, but their model-dependent proposals and realized point sets
need not remain identical after training diverges. The
corrected-versus-uncorrected contrast concerns the full paired training
trajectories.

\hypertarget{problems}{%
\subsection{Problems}\label{problems}}

The benchmark suite contains viscous Burgers, a manufactured Allen--Cahn
problem, a standing-wave equation and a manufactured two-dimensional
Helmholtz problem. An initial high-frequency Helmholtz screen was not
used for method ranking because every method had relative error above
one. We calibrated a solvable configuration using only M0 and a
predetermined criterion \(L_2<0.2\), then froze it before comparing
adaptive methods.

Four Burgers observation regimes assess robustness: thirteen clean
observations; thirteen observations with Gaussian noise and outliers;
one observation with strong heteroscedastic noise; and a missing sensor
block with correlated noise.

The governing equations and exact or independent targets are:

\begin{itemize}
\tightlist
\item
  \textbf{Burgers:} \(u_t+u u_x-\nu u_{xx}=0\), with \(\nu=0.01/\pi\),
  \(x\in[-1,1]\), \(t\in[0,1]\), \(u(0,x)=-\sin(\pi x)\) and zero
  Dirichlet boundaries. V10 uses the Cole--Hopf integral solution
  (Burkardt n.d.).
\item
  \textbf{Manufactured Allen--Cahn:} \(u_t-Du_{xx}-\rho(u-u^3)=f\), with
  \(D=10^{-3}\), \(\rho=5\), the same space-time domain and
  \(u=e^{-t}\sin(\pi x)\). The forcing is \(f=-u+D\pi^2u-\rho(u-u^3)\);
  exact initial and boundary values are imposed softly. This is not the
  difficult unforced benchmark.
\item
  \textbf{Wave:} \(u_{tt}-c^2u_{xx}=0\), with \(c=1\) on \([0,1]^2\),
  \(u=\sin(\pi x)\cos(\pi t)\), zero boundary values, initial
  displacement \(\sin(\pi x)\) and zero initial velocity. Displacement
  and velocity MSEs are summed in the initial-condition component.
\item
  \textbf{Manufactured Helmholtz:} \(\Delta u+k^2u=f\) on \([0,1]^2\),
  with \(k=1\), \(u=\sin(\pi x)\sin(\pi y)\) and \(f=(1-2\pi^2)u\), with
  zero boundaries.
\end{itemize}

All final V09--V10 evaluations concern these four PDE families. The
damped oscillator mentioned in the original research blueprint is not
included in the reported experimental suite. All physical parameters
above are prescribed; the present experiments do not identify unknown
PDE coefficients from data.

\hypertarget{experimental-controls}{%
\subsection{Experimental controls}\label{experimental-controls}}

The historical ten-seed comparisons and the exploratory V09--V10
campaigns use ten paired seeds
\texttt{11,\ 22,\ 33,\ 44,\ 55,\ 66,\ 77,\ 88,\ 99,\ 111}. Methods share
architecture, initialization seed, observation realization, collocation
cardinality and clean evaluation target. All configurations and
environment metadata are stored with each run. Failed runs are retained
and reported rather than silently discarded.

The primary metric is relative \(L_2\) error. Secondary metrics include
maximum absolute error, component-weight total variation, the
total-variation distance between component-gradient distributions on
training and audit points, the logarithmic physics-gradient gap, wall
time, residual evaluations, importance- weight effective sample size,
and KDE pairwise-distance counts.

We report medians, quartiles, paired median differences, 95\% bootstrap
intervals and two-sided sign tests. A method is not called superior when
only its point estimate improves or when the paired interval crosses
zero.

\hypertarget{v09v10-extensions-and-common-loss-audit}{%
\subsection{V09--V10 extensions and common-loss
audit}\label{v09v10-extensions-and-common-loss-audit}}

V09 consists of 200 runs: four PDE configurations, five
local/global/potential variants and ten seeds. V10 adds 200 Burgers
observation-regime runs, 200 doubled-step runs and 160 common-loss
baseline runs. The latter reuse 40 existing V10 vRBA checkpoints for
comparison. Each suite freezes its manifest before training but follows
earlier empirical findings; these suites are exploratory, not
independent confirmations of hypotheses selected using V07 or V09.

The four factorial arms are neither, global\_only, local\_only and
full\_exp; full\_quad changes only the potential relative to full\_exp.
V10 uses MSE for observations in every arm. In historical noisy-data
comparisons, local vRBA multipliers triggered a squared residual loss
while several non-vRBA methods used Huber. Consequently, those
comparisons contrast complete recipes and do not identify an effect of
attention alone. V10 synthetic Burgers observations and evaluation
targets use the independent Cole--Hopf solution.

For each suite, four predeclared contrasts per case form a family of 16
tests for relative L2, with a separate family for maximum absolute
error. We use two-sided exact sign tests, Holm correction, and
10000-resample percentile bootstrap intervals for the paired median
difference. Intervals are marginal, not simultaneous. Initialization and
observation tensors are fingerprinted; the 400 factorial/budget V10
checkpoints are re-evaluated independently, including a denser grid; the
160 additional baseline checkpoints are independently checked on the
original scoring grid.

\hypertarget{architecture-budgets-and-observation-counts}{%
\subsection{Architecture, budgets and observation
counts}\label{architecture-budgets-and-observation-counts}}

All networks use tanh activations, Adam with learning rate \(10^{-3}\),
float64 arithmetic, gradient clipping at 100 and one numerical CPU
thread per process. Table 1 summarizes the four base configurations.
Architectures and point cardinalities are matched between methods within
each configuration, not globally across all PDEs.

\begingroup\small

\begin{longtable}[]{@{}
  >{\raggedright\arraybackslash}p{(\columnwidth - 8\tabcolsep) * \real{0.1765}}
  >{\raggedright\arraybackslash}p{(\columnwidth - 8\tabcolsep) * \real{0.1765}}
  >{\raggedright\arraybackslash}p{(\columnwidth - 8\tabcolsep) * \real{0.1765}}
  >{\raggedleft\arraybackslash}p{(\columnwidth - 8\tabcolsep) * \real{0.2353}}
  >{\raggedleft\arraybackslash}p{(\columnwidth - 8\tabcolsep) * \real{0.2353}}@{}}
\caption{Base architectures and optimization budgets.}\tabularnewline
\toprule\noalign{}
\begin{minipage}[b]{\linewidth}\raggedright
Configuration
\end{minipage} & \begin{minipage}[b]{\linewidth}\raggedright
Hidden layers x width
\end{minipage} & \begin{minipage}[b]{\linewidth}\raggedright
Physics / initial / boundary points
\end{minipage} & \begin{minipage}[b]{\linewidth}\raggedleft
V09 steps
\end{minipage} & \begin{minipage}[b]{\linewidth}\raggedleft
V10 doubled steps
\end{minipage} \\
\midrule\noalign{}
\endfirsthead
\toprule\noalign{}
\begin{minipage}[b]{\linewidth}\raggedright
Configuration
\end{minipage} & \begin{minipage}[b]{\linewidth}\raggedright
Hidden layers x width
\end{minipage} & \begin{minipage}[b]{\linewidth}\raggedright
Physics / initial / boundary points
\end{minipage} & \begin{minipage}[b]{\linewidth}\raggedleft
V09 steps
\end{minipage} & \begin{minipage}[b]{\linewidth}\raggedleft
V10 doubled steps
\end{minipage} \\
\midrule\noalign{}
\endhead
\bottomrule\noalign{}
\endlastfoot
Burgers long & 4 x 48 & 512 / 128 / 128 & 2500 & 5000 \\
Allen--Cahn & 3 x 32 & 256 / 64 / 64 & 600 & 1200 \\
Helmholtz & 4 x 48 & 512 / 0 / 128 & 1500 & 3000 \\
Wave & 3 x 32 & 256 / 64 / 64 & 600 & 1200 \\
\end{longtable}

\endgroup

The four difficult-observation Burgers cases use 3 x 32 networks, 256
physics, 64 initial and 64 boundary points, and 600 steps. Their
realized observation counts are 13 (sparse clean), 13 (sparse noisy), 1
(extreme) and 64 (missing block). The common-loss extension uses MSE
throughout. Detailed noise settings and all paired contrasts are
supplied in the supplement.

\hypertarget{results}{%
\section{Results}\label{results}}

Sections 4.1--4.6 report historical configurations and references. New
V09--V10 evidence follows in Sections 4.7--4.11. Historical noisy-data
rankings are recipe comparisons subject to the loss confound described
in Section 3.5.

\hypertarget{fixed-reference-decoupling}{%
\subsection{Fixed-reference
decoupling}\label{fixed-reference-decoupling}}

On long Burgers runs, M5 shows a positive concentration--gradient-gap
correlation in all ten seeds. Nevertheless, M6 does not show a
conclusive improvement in accuracy under the reported paired interval:
the paired median difference M6−M5 is \(-0.0291\), with 95\% interval
\([-0.1208,0.1678]\). Nor does it consistently reduce weight variation.
Because controller cardinality and residual-evaluation counts are
matched, this null result cannot be attributed to additional reference
points.

\hypertarget{exact-correction-and-its-matched-control}{%
\subsection{Exact correction and its matched
control}\label{exact-correction-and-its-matched-control}}

M7 substantially reduces the mean physics-gradient log-gap relative to
M7U on Burgers: paired median difference \(-0.4351\), 95\% interval
\([-0.6182,-0.2520]\), favorable in all ten seeds. Across the four
difficult observation regimes, M7 also improves relative \(L_2\) error
over M7U by \(-0.0483\), \(-0.0482\), \(-0.0229\), and \(-0.0360\),
respectively; every bootstrap interval excludes zero.

This effect is not universal. On Allen--Cahn, M7 is worse than M7U by a
paired median of \(+0.0097\), with a 95\% interval just above zero.
Thus, restoring the chosen target measure may improve diagnostic
alignment while worsening the optimization trajectory for a particular
PDE.

\hypertarget{comparison-with-kde-volume-weighting}{%
\subsection{Comparison with KDE volume
weighting}\label{comparison-with-kde-volume-weighting}}

On the standing-wave problem, fixed-weight VW is better than M7 in all
ten seeds: M7−VW is \(+0.0280\), 95\% interval \([0.0199,0.0483]\). On
calibrated Helmholtz, the ranking reverses: M7−VW is \(-0.0181\), 95\%
interval \([-0.0235,-0.0112]\), favorable to M7 in all ten seeds.
Burgers long and Allen--Cahn show no conclusive difference between M7
and VW.

At the default KDE bandwidth, M7 outperforms VW-CA in each of the four
difficult Burgers observation regimes, with paired intervals excluding
zero. However, M7 does not outperform the simpler fixed-weight VW in any
of those four regimes. The addition of a global adaptive controller
therefore does not automatically improve a locally corrected loss.

\hypertarget{bandwidth-sensitivity}{%
\subsection{Bandwidth sensitivity}\label{bandwidth-sensitivity}}

The KDE choice changes the conclusion. On noisy Burgers, reducing the
bandwidth to 0.5 times Scott's rule yields median errors 0.4742 for VW
and 0.5527 for M7. The paired M7−VW difference is \(+0.0710\), 95\%
interval \([0.0161,0.1165]\). At that bandwidth, the previously
favorable M7−VW-CA contrast also becomes inconclusive. On the wave
equation, VW remains better than M7 at all three tested bandwidth
scales.

\hypertarget{computational-cost}{%
\subsection{Computational cost}\label{computational-cost}}

In a dedicated sequential CPU profile, M7 adds a median 10.5\% wall-time
cost relative to M5, 7.2\% relative to M6 and 4.9\% relative to M7U,
while retaining the same number of residual evaluations. For 256
collocation points and 600 steps, the dense VW correction additionally
evaluates 1,638,400 pairwise KDE distances per run. The correction part
of M7 is linear after the proposal is known, whereas dense KDE is
quadratic in the number of collocation points. This theoretical
advantage has not yet been verified by a dedicated multi-size CPU/GPU
profile.

\hypertarget{variational-residual-based-attention}{%
\subsection{Variational residual-based
attention}\label{variational-residual-based-attention}}

The protocol-frozen common-backbone vRBA baseline completes all runs
without non-finite values. Its final attention effective-sample fraction
ranges from 0.785 to 0.990, confirming non-uniform but non-collapsed
local weighting.

vRBA obtains median relative L2 errors of 0.2507 on Burgers, 0.0401 on
Allen--Cahn, 0.0388 on calibrated Helmholtz, and 0.1745 on the wave
equation. Against M7, paired median differences are respectively -0.1770
{[}-0.2339, -0.0114{]}, -0.0648 {[}-0.0903, -0.0509{]}, -0.0611
{[}-0.0690, -0.0496{]}, and -0.1322 {[}-0.2163, -0.1190{]}. Against
fixed-weight VW, they are -0.1236 {[}-0.2566, -0.0283{]}, -0.0617
{[}-0.1269, -0.0267{]}, -0.0844 {[}-0.0891, -0.0608{]}, and -0.1061
{[}-0.1502, -0.0836{]}.

The same L2 signal persists across clean sparse data, noisy and outlier-
contaminated data, one strongly corrupted observation, and a missing
sensor block. However, the extreme one-observation regime has only eight
favorable seeds out of ten against M7 and VW, yielding a two-sided
sign-test p-value of 0.109 despite bootstrap intervals below zero.

The maximum-error result contradicts a universal robustness
interpretation. On long Burgers runs, vRBA's median maximum error is
1.7209, compared with 1.3429 for M5 and 1.2993 for M6. The paired
vRBA-M5 difference is +0.3336 {[}0.0265, 0.5694{]}, and vRBA-M6 is
+0.4011 {[}0.2067, 0.5759{]}. Thus, vRBA improves the global L2 norm
while worsening the largest localized error. Observed runtime overhead
relative to M5 ranges from 10.4\% to 81.8\% across PDEs.

Spatial analysis resolves this apparent contradiction. The median across
seeds of the gridwise 95th absolute-error percentile is 0.1027 for vRBA,
versus 0.6698 for M5 and 0.6907 for M6; the corresponding 99th
percentiles are 0.8588, 1.1470 and 1.0516. vRBA therefore improves
almost the entire domain while leaving a very narrow late-time error
spike near the shock. In the median error field, this spike peaks around
\(t=0.975\), \(x=-0.025\).

\hypertarget{localglobal-decomposition-at-the-v09-budgets}{%
\subsection{Local/global decomposition at the V09
budgets}\label{localglobal-decomposition-at-the-v09-budgets}}

The 200-run V09 campaign found that full exponential vRBA improves
relative L2 against both local-only and global-only variants for the
calibrated Helmholtz case (k=1). Both contrasts and the additive-scale
interaction pass Holm at p=0.03125. For the wave case, the quadratic
potential improves L2 against the exponential potential (Holm
p=0.03125). No planned Burgers contrast passes the multiplicity
correction. These are configuration-specific findings, not evidence of a
universally optimal potential or controller.

V09 separately measured 60 CPU windows using three seeds, five variants
and four PDEs. Fifty-nine profiling trials completed; the last trial
ended after the prespecified timing window. The full exponential variant
has median paired per-step cost ratios of about 1.71--1.86 relative to
neither. These are local per-step measurements, not time-to-accuracy or
equal-time rankings.

\hypertarget{independent-burgers-reference}{%
\subsection{Independent Burgers
reference}\label{independent-burgers-reference}}

Cole--Hopf quadrature at orders 128 and 256 agrees within 9.11e-15 on
the tested 161-by-321 grid. Adaptive quadrature independently verifies
20 selected points within 4.44e-16. The historical Rusanov reference
with 257 spatial points differs from this solution by 0.03410 in
relative L2 and 0.28578 in maximum error; at 1025 points these
discrepancies fall to 0.00733 and 0.06969.

Re-evaluating all 50 frozen V09 Burgers models on the dense grid with
Cole--Hopf gives median L2 values 0.30724, 0.26523, 0.22496, 0.26369 and
0.26370 for neither, global\_only, local\_only, full\_exp and full\_quad
respectively. No planned contrast passes Holm in this exploratory
four-contrast reanalysis. The numerical-reference discrepancy does not
explain away the observed large model errors, but it prevents treating
historical maximum errors as exact.

\hypertarget{sparse-and-corrupted-observations-under-a-common-loss}{%
\subsection{Sparse and corrupted observations under a common
loss}\label{sparse-and-corrupted-observations-under-a-common-loss}}

All 200 runs finish successfully and all 40 five-arm initialization
blocks match. The table reports median relative L2 on clean Cole--Hopf
targets.

\begingroup\small

\begin{longtable}[]{@{}lrrrrr@{}}
\caption{Median relative L2 error under difficult
observations.}\tabularnewline
\toprule\noalign{}
Case & Neither & Global only & Local only & Full exp. & Full quad. \\
\midrule\noalign{}
\endfirsthead
\toprule\noalign{}
Case & Neither & Global only & Local only & Full exp. & Full quad. \\
\midrule\noalign{}
\endhead
\bottomrule\noalign{}
\endlastfoot
extreme & 0.53812 & 0.53812 & 0.44694 & 0.45883 & 0.47695 \\
sparse\_noisy & 0.46306 & 0.46373 & 0.37333 & 0.37603 & 0.41816 \\
sparse\_clean & 0.46061 & 0.45779 & 0.38439 & 0.36866 & 0.40329 \\
block\_missing & 0.44003 & 0.43793 & 0.34951 & 0.34398 & 0.34288 \\
\end{longtable}

\endgroup

After correction, full\_exp improves L2 over global\_only for the
missing-block case (paired median difference -0.06999, marginal 95\% CI
{[}-0.09386,-0.03466{]}, Holm p=0.03125). For extreme data, full\_exp
improves the maximum error over full\_quad (difference -0.06696, CI
{[}-0.12910,-0.00587{]}, Holm p=0.03125). Other planned contrasts do not
pass Holm. Section 4.11 separately retests the historical M5/M6/M7/VW
comparisons with a common observation loss.

\hypertarget{doubled-optimization-budgets}{%
\subsection{Doubled optimization
budgets}\label{doubled-optimization-budgets}}

The 200 additional runs use 5000 Burgers, 1200 Allen--Cahn, 3000
Helmholtz and 1200 wave steps. The following median relative L2 values
come from independent checkpoint evaluations. Initial tensors match both
within each five-arm block and against the historical V09 counterparts.
No model selection is based on the evaluation scores.

\begingroup\small

\begin{longtable}[]{@{}lrrrrr@{}}
\caption{Median relative L2 error at doubled optimization
budgets.}\tabularnewline
\toprule\noalign{}
Case & Neither & Global only & Local only & Full exp. & Full quad. \\
\midrule\noalign{}
\endfirsthead
\toprule\noalign{}
Case & Neither & Global only & Local only & Full exp. & Full quad. \\
\midrule\noalign{}
\endhead
\bottomrule\noalign{}
\endlastfoot
helmholtz & 0.08717 & 0.01917 & 0.03245 & 0.00523 & 0.02076 \\
allen\_cahn & 0.02859 & 0.02873 & 0.02573 & 0.02785 & 0.02291 \\
burgers & 0.30754 & 0.24549 & 0.20305 & 0.25680 & 0.26014 \\
wave & 0.12374 & 0.13976 & 0.02578 & 0.02747 & 0.02652 \\
\end{longtable}

\endgroup

Planned contrasts passing Holm in the doubled-step suite are listed
below; the complete table, including null and adverse contrasts,
accompanies the code.

\begingroup\small

\begin{longtable}[]{@{}
  >{\raggedright\arraybackslash}p{(\columnwidth - 10\tabcolsep) * \real{0.1500}}
  >{\raggedright\arraybackslash}p{(\columnwidth - 10\tabcolsep) * \real{0.1500}}
  >{\raggedright\arraybackslash}p{(\columnwidth - 10\tabcolsep) * \real{0.1500}}
  >{\raggedleft\arraybackslash}p{(\columnwidth - 10\tabcolsep) * \real{0.2000}}
  >{\raggedright\arraybackslash}p{(\columnwidth - 10\tabcolsep) * \real{0.1500}}
  >{\raggedleft\arraybackslash}p{(\columnwidth - 10\tabcolsep) * \real{0.2000}}@{}}
\caption{Planned doubled-budget contrasts passing Holm
correction.}\tabularnewline
\toprule\noalign{}
\begin{minipage}[b]{\linewidth}\raggedright
Case
\end{minipage} & \begin{minipage}[b]{\linewidth}\raggedright
Metric
\end{minipage} & \begin{minipage}[b]{\linewidth}\raggedright
Contrast
\end{minipage} & \begin{minipage}[b]{\linewidth}\raggedleft
Paired median difference
\end{minipage} & \begin{minipage}[b]{\linewidth}\raggedright
Marginal 95\% CI
\end{minipage} & \begin{minipage}[b]{\linewidth}\raggedleft
Holm p
\end{minipage} \\
\midrule\noalign{}
\endfirsthead
\toprule\noalign{}
\begin{minipage}[b]{\linewidth}\raggedright
Case
\end{minipage} & \begin{minipage}[b]{\linewidth}\raggedright
Metric
\end{minipage} & \begin{minipage}[b]{\linewidth}\raggedright
Contrast
\end{minipage} & \begin{minipage}[b]{\linewidth}\raggedleft
Paired median difference
\end{minipage} & \begin{minipage}[b]{\linewidth}\raggedright
Marginal 95\% CI
\end{minipage} & \begin{minipage}[b]{\linewidth}\raggedleft
Holm p
\end{minipage} \\
\midrule\noalign{}
\endhead
\bottomrule\noalign{}
\endlastfoot
wave & relative\_l2 & full\_exp-global\_only & -0.10001 & {[}-0.14629,
-0.08958{]} & 0.03125 \\
helmholtz & max\_abs\_error & interaction & +0.11349 & {[}+0.08393,
+0.13949{]} & 0.03125 \\
wave & max\_abs\_error & full\_exp-global\_only & -0.19593 &
{[}-0.30415, -0.14881{]} & 0.03125 \\
\end{longtable}

\endgroup

The quadratic-potential L2 advantage established for the V09 wave
configuration does not pass Holm at the doubled horizon. Full\_exp still
improves both L2 and maximum error over global\_only for waves. A
positive Helmholtz maximum-error interaction reflects the additive
contrast and does not imply that full\_exp is the least accurate model.

The doubled-step comparison measures sensitivity to optimization
duration. It is not an equal-compute or equal-time experiment, and
historical and V10 wall times must not be pooled as a controlled timing
comparison.

\hypertarget{re-running-the-historical-baselines-with-a-common-observation-loss}{%
\subsection{Re-running the historical baselines with a common
observation
loss}\label{re-running-the-historical-baselines-with-a-common-observation-loss}}

An additional 160 runs cover M5, M6, M7 and VW across the same four
observation regimes and ten seeds, using the exact V10 MSE/Cole--Hopf
configurations. Forty full\_exp checkpoints from Section 4.9 are reused
as vRBA comparators; they are not new or independent training runs.
Initial parameter and observation tensors match across all five compared
methods. The table gives median L2.

\begingroup\small

\begin{longtable}[]{@{}lrrrrr@{}}
\caption{Common-loss baseline median relative L2 errors.}\tabularnewline
\toprule\noalign{}
Case & M5 & M6 & M7 & VW & vRBA \\
\midrule\noalign{}
\endfirsthead
\toprule\noalign{}
Case & M5 & M6 & M7 & VW & vRBA \\
\midrule\noalign{}
\endhead
\bottomrule\noalign{}
\endlastfoot
block\_missing & 0.56166 & 0.55393 & 0.53397 & 0.48595 & 0.34398 \\
sparse\_clean & 0.58673 & 0.55823 & 0.55471 & 0.49238 & 0.36866 \\
sparse\_noisy & 0.57248 & 0.55808 & 0.54189 & 0.50375 & 0.37603 \\
extreme & 0.61185 & 0.60330 & 0.59146 & 0.56037 & 0.45883 \\
\end{longtable}

\endgroup

All newly trained checkpoints have been independently re-evaluated. The
following contrasts pass Holm in the separately declared family of 16
per metric. All contrasts, including null and unfavorable effects, are
retained in the accompanying paired\_statistics.csv.

\begingroup\small

\begin{longtable}[]{@{}
  >{\raggedright\arraybackslash}p{(\columnwidth - 10\tabcolsep) * \real{0.1500}}
  >{\raggedright\arraybackslash}p{(\columnwidth - 10\tabcolsep) * \real{0.1500}}
  >{\raggedright\arraybackslash}p{(\columnwidth - 10\tabcolsep) * \real{0.1500}}
  >{\raggedleft\arraybackslash}p{(\columnwidth - 10\tabcolsep) * \real{0.2000}}
  >{\raggedright\arraybackslash}p{(\columnwidth - 10\tabcolsep) * \real{0.1500}}
  >{\raggedleft\arraybackslash}p{(\columnwidth - 10\tabcolsep) * \real{0.2000}}@{}}
\caption{Common-loss paired contrasts passing Holm
correction.}\tabularnewline
\toprule\noalign{}
\begin{minipage}[b]{\linewidth}\raggedright
Case
\end{minipage} & \begin{minipage}[b]{\linewidth}\raggedright
Metric
\end{minipage} & \begin{minipage}[b]{\linewidth}\raggedright
Contrast
\end{minipage} & \begin{minipage}[b]{\linewidth}\raggedleft
Paired median difference
\end{minipage} & \begin{minipage}[b]{\linewidth}\raggedright
Marginal 95\% CI
\end{minipage} & \begin{minipage}[b]{\linewidth}\raggedleft
Holm p
\end{minipage} \\
\midrule\noalign{}
\endfirsthead
\toprule\noalign{}
\begin{minipage}[b]{\linewidth}\raggedright
Case
\end{minipage} & \begin{minipage}[b]{\linewidth}\raggedright
Metric
\end{minipage} & \begin{minipage}[b]{\linewidth}\raggedright
Contrast
\end{minipage} & \begin{minipage}[b]{\linewidth}\raggedleft
Paired median difference
\end{minipage} & \begin{minipage}[b]{\linewidth}\raggedright
Marginal 95\% CI
\end{minipage} & \begin{minipage}[b]{\linewidth}\raggedleft
Holm p
\end{minipage} \\
\midrule\noalign{}
\endhead
\bottomrule\noalign{}
\endlastfoot
block\_missing & relative\_l2 & vrba-m5 & -0.19580 & {[}-0.24059,
-0.18515{]} & 0.03125 \\
block\_missing & relative\_l2 & vrba-m6 & -0.19318 & {[}-0.24042,
-0.18069{]} & 0.03125 \\
block\_missing & relative\_l2 & vrba-m7 & -0.17172 & {[}-0.20840,
-0.12736{]} & 0.03125 \\
block\_missing & relative\_l2 & vrba-vw & -0.13158 & {[}-0.15465,
-0.09959{]} & 0.03125 \\
sparse\_clean & relative\_l2 & vrba-m5 & -0.22439 & {[}-0.32713,
-0.15273{]} & 0.03125 \\
sparse\_clean & relative\_l2 & vrba-m6 & -0.21344 & {[}-0.32406,
-0.15757{]} & 0.03125 \\
sparse\_clean & relative\_l2 & vrba-m7 & -0.17495 & {[}-0.24497,
-0.13279{]} & 0.03125 \\
sparse\_noisy & relative\_l2 & vrba-m5 & -0.19789 & {[}-0.30269,
-0.17423{]} & 0.03125 \\
extreme & relative\_l2 & vrba-m5 & -0.13299 & {[}-0.20818, -0.06938{]} &
0.03125 \\
extreme & relative\_l2 & vrba-m6 & -0.13806 & {[}-0.20866, -0.11213{]} &
0.03125 \\
extreme & relative\_l2 & vrba-m7 & -0.13349 & {[}-0.15551, -0.06411{]} &
0.03125 \\
\end{longtable}

\endgroup

vRBA has the lowest median L2 in all four regimes, with 11 of 16 L2
contrasts passing Holm; no maximum-error contrast passes its separate
correction. This targeted rerun removes the observation-loss mismatch in
these cases. Its protocol followed the V10 audit and remains
exploratory. It does not establish performance on unseen physical
systems or a general robustness guarantee.

The following figures are generated from the stored numerical analyses.
Figure 1 documents the Burgers reference discrepancy; Figure 2 displays
every paired seed in the common-loss comparison; Figure 3 shows
doubled-step ablations. They must be interpreted together with the
corrected statistics.

\begin{figure}
\centering
\includegraphics[width=0.95\textwidth,height=\textheight]{outputs/reference_v10_analysis/reference_validation.png}
\caption{Burgers reference validation and re-evaluation of frozen V09
models. The second panel also changes the evaluation grid; it does not
isolate the reference change alone.}
\end{figure}

\begin{figure}
\centering
\includegraphics[width=0.8\textwidth,height=\textheight]{submission/v11/common_loss_by_seed.pdf}
\caption{Common-MSE Burgers comparisons. Points represent the ten
individual seeds and black bars show medians. No inference is based
solely on these visual rankings.}
\end{figure}

\begin{figure}
\centering
\includegraphics[width=0.8\textwidth,height=\textheight]{outputs/budget_v10_analysis/l2_by_seed.png}
\caption{Doubled-step ablations across four PDE configurations. Budgets
match optimization steps within each PDE and do not equalize wall time.}
\end{figure}

\hypertarget{discussion}{%
\section{Discussion}\label{discussion}}

The experiments evaluate the original fixed-reference hypothesis and
separate three statements that are often conflated. First, adaptive
collocation can change what a loss controller observes. Second, removing
this change can alter gradients. Third, neither fact guarantees lower
solution error. M6 demonstrates that an independent reference set can
change the controller without producing a robust benefit. M7
demonstrates that exact proposal correction can restore a target
finite-population mean and improve a matched control in data-poor
Burgers regimes. VW demonstrates that an estimated geometric correction
can be equally or more effective, especially when its bandwidth is
favorable and global component weights remain fixed.

The poor behavior of VW-CA relative to VW in several cases suggests an
interaction between local volume weights and global component balancing.
One possible explanation is that squared local volume weights change the
magnitude and variance of the physics gradient, after which the global
controller reacts to a signal that is already strongly rescaled. This
interpretation is an inference, not yet a demonstrated mechanism. A
factorial experiment varying local correction, global controller and
update frequency is needed.

The vRBA comparison sharpens the distinction between restoring a target
measure and deliberately changing the training objective. In this
benchmark, the tilted objective is much more effective in global L2
error than either exact or KDE-based correction. Yet the Burgers
maximum-error result shows that this benefit is not uniform over the
domain. The spatial maps document localized error patterns alongside
lower global error. They do not establish that attention causes the
localized spike. The V09--V10 ablations vary local/global mechanisms
under fixed point sets, while the historical recipe comparisons also
differ in sampling and controller details.

The practical conclusion is not that one correction should replace every
other. When the proposal probability is exactly available, M7 avoids
density estimation and has favorable asymptotic correction cost. When
only a realized nonuniform point cloud is available, volume weighting
remains applicable. If the objective is deliberately tilted toward high
residuals rather than restored to a uniform measure, a variational
adaptive method such as vRBA may be more appropriate.

\hypertarget{limitations}{%
\section{Limitations}\label{limitations}}

\begin{enumerate}
\def\labelenumi{\arabic{enumi}.}
\tightlist
\item
  The VW implementation reproduces its core weighting mechanism on a
  common backbone, not the complete published architecture and
  Adam--L-BFGS schedule.
\item
  vRBA is implemented on the common backbone, but the complete official
  JAX architecture, hard constraints and long optimizer schedule have
  not been reproduced.
\item
  The original high-frequency Helmholtz case remains unsolved; the
  calibrated case answers a different, easier question and is reported
  separately.
\item
  Ten seeds provide useful paired evidence but limited power for small
  effects.
\item
  KDE bandwidth sensitivity was tested on two problems, not the full
  suite.
\item
  Timing conclusions are CPU-specific, and the KDE scaling claim still
  needs controlled multi-size and GPU measurements.
\item
  The PDE suite is synthetic. No conclusion currently transfers to
  tokamak plasma data or real-time instability control.
\item
  Local/global and potential ablations, spatial maps and doubled-step
  sensitivity are available. Common-loss noisy baseline reruns are also
  available. Equal-time comparisons remain outside the claims of this
  paper. A separate-host portability protocol is supplied; its execution
  status is reported in the supplement.
\item
  The public literature review is reproducible but not exhaustive across
  subscription databases.
\end{enumerate}

\hypertarget{conclusion}{%
\section{Conclusion}\label{conclusion}}

Residual sampling and adaptive loss balancing form a feedback system
whose behavior depends on the measure used to estimate each loss
component. The tested fixed-reference ReCoA-PINN implementation does not
establish general stabilization across the evaluated configurations.
Exact proposal reweighting can reduce distribution-induced gradient
discrepancies and improve matched controls in sparse/noisy Burgers
regimes, but it does not dominate KDE volume weighting across PDEs or
bandwidths. The common-backbone vRBA recipe shows strong historical
relative-L2 performance, qualified by a noisy-data loss confound and
substantial maximum-error and cost tradeoffs. The new common-loss
ablations and doubled budgets identify conditional rather than universal
benefits. Robust claims about adaptive PINNs therefore require matched
controls, several PDEs and seeds, multiple error metrics, sensitivity
analyses, and explicit separation between restoring a measure and
optimizing a deliberately tilted objective.

\hypertarget{data-code-and-reproducibility-statement}{%
\section{Data, code and reproducibility
statement}\label{data-code-and-reproducibility-statement}}

The PyTorch implementation stores resolved YAML configurations, software
and hardware metadata, metric trajectories, model checkpoints and final
summaries for every run. Campaign manifests regenerate paired tables and
figures while resuming only missing cells. The repository contains
automated unit and smoke tests for samplers, corrections, PDE residuals,
reproducibility and reporting. V09 supplies 200 ablation runs and V10
adds 560 verified runs, in addition to the historical baseline
campaigns. The V10 numerical suite passed 55 tests. Raw campaign counts
distinguish completed training, repeated profiling and reused historical
evidence; they must not be summed as independent physical experiments.

\hypertarget{independent-host-verification}{%
\section{Independent-host
verification}\label{independent-host-verification}}

On 6 September 2026, a separate GitHub-hosted Ubuntu runner
reconstructed the scores of all 560 frozen V10 models and all 96 paired
contrasts within the predeclared tolerance (relative \(10^{-8}\),
absolute \(10^{-10}\)). The 55 automated tests passed. All 12
preselected retrainings completed and their L2 and maximum errors met
the predeclared numerical tolerance (relative \(10^{-3}\), absolute
\(10^{-5}\)). The largest absolute differences were
\(9.53\times10^{-9}\) in L2 and \(1.40\times10^{-8}\) in maximum error.

The stricter joint gate nevertheless failed: initial model fingerprints
matched in all 12 cases, but training/reference-set fingerprints
differed in the four noisy-observation cases. Only 8/12 satisfied every
declared condition. The environments differ in host and PyTorch version
(2.14.0+cpu versus 2.8.0+cpu). The original raw initial tensors were not
retained, so the differing bytes cannot be located retrospectively. A
floating-point or library-version effect is plausible but not
established by these hashes. This is successful numerical reconstruction
with a documented input-identity limitation, not bitwise reproduction of
the complete experiment. The failed gate and its thresholds are
preserved in the repository and supplement.

\hypertarget{author-metadata-and-declarations}{%
\section{Author metadata and
declarations}\label{author-metadata-and-declarations}}

This is an author-review copy. Youssef EL MOUTEE is the identified
working author. The final author list and order, affiliations,
corresponding-author details, contributions, funding and
competing-interest statements require author confirmation before
submission. No unconfirmed affiliation or consent has been assigned to
the supervisor. The code and archived evidence are in the public
repository \href{https://github.com/ucef90/Doctorat}{ucef90/Doctorat},
with frozen V10 evidence at commit
\texttt{4161ed57f232ea7b796480b38201502b3cabbf67}. All journal-specific
declarations require author confirmation before submission.

\hypertarget{references}{%
\section*{References}\label{references}}
\addcontentsline{toc}{section}{References}

\hypertarget{refs}{}
\begin{CSLReferences}{1}{0}
\leavevmode\vadjust pre{\hypertarget{ref-burkardt}{}}%
Burkardt, John. n.d. {``Burgers\_exact: Integral Solution of the Burgers
Equation.''} Accessed September 6, 2026.
\url{https://people.sc.fsu.edu/~jburkardt/f77_src/burgers_exact/burgers_exact.html}.

\leavevmode\vadjust pre{\hypertarget{ref-celaya2025}{}}%
Celaya, Adrian, David Fuentes, and Beatrice Riviere. 2025. {``Adaptive
Collocation Point Strategies For Physics Informed Neural Networks via
the QR Discrete Empirical Interpolation Method.''} arXiv preprint.
\url{https://arxiv.org/abs/2501.07700v4}.

\leavevmode\vadjust pre{\hypertarget{ref-chen2025}{}}%
Chen, Wenqian, Amanda Howard, and Panos Stinis. 2025. {``Self-Adaptive
Weighting and Sampling for Physics-Informed Neural Networks.''} arXiv
preprint. \url{https://arxiv.org/abs/2511.05452v2}.

\leavevmode\vadjust pre{\hypertarget{ref-lau2024}{}}%
Lau, Gregory Kang Ruey, Apivich Hemachandra, See-Kiong Ng, and Bryan
Kian Hsiang Low. 2024. {``PINNACLE: PINN Adaptive ColLocation and
Experimental Points Selection.''} ICLR 2024; arXiv version.
\url{https://arxiv.org/abs/2404.07662}.

\leavevmode\vadjust pre{\hypertarget{ref-nabian2021}{}}%
Nabian, Mohammad Amin, Rini Jasmine Gladstone, and Hadi Meidani. 2021.
{``Efficient Training of Physics‐informed Neural Networks via Importance
Sampling.''} \emph{Computer-Aided Civil and Infrastructure Engineering}
36 (August): 962--77. \url{https://doi.org/10.1111/mice.12685}.

\leavevmode\vadjust pre{\hypertarget{ref-raissi2019}{}}%
Raissi, M., P. Perdikaris, and G.E. Karniadakis. 2019.
{``Physics-Informed Neural Networks: A Deep Learning Framework for
Solving Forward and Inverse Problems Involving Nonlinear Partial
Differential Equations.''} \emph{Journal of Computational Physics} 378
(February): 686--707. \url{https://doi.org/10.1016/j.jcp.2018.10.045}.

\leavevmode\vadjust pre{\hypertarget{ref-singh2026}{}}%
Singh, Divyavardhan, Shubham Kamble, Dimple Sonone, and Kishor Upla.
2026. {``Stabilized Adaptive Loss and Residual-Based Collocation for
Physics-Informed Neural Networks.''} arXiv preprint.
\url{https://arxiv.org/abs/2603.03224}.

\leavevmode\vadjust pre{\hypertarget{ref-song2025}{}}%
Song, Jiahao, Wenbo Cao, Fei Liao, and Weiwei Zhang. 2025. {``VW-PINNs:
A Volume Weighting Method for PDE Residuals in Physics-Informed Neural
Networks.''} \emph{Acta Mechanica Sinica} 41 (March): 324140.
\url{https://doi.org/10.1007/s10409-024-24140-x}.

\leavevmode\vadjust pre{\hypertarget{ref-toscano2026}{}}%
Toscano, Juan Diego, Daniel T. Chen, Vivek Ooomen, Jérôme Darbon, and
George Em Karniadakis. 2026. {``A Variational Framework for
Residual-Based Adaptivity in Neural PDE Solvers and Operator
Learning.''} \emph{Npj Artificial Intelligence} 2 (March): 32.
\url{https://doi.org/10.1038/s44387-026-00084-4}.

\leavevmode\vadjust pre{\hypertarget{ref-visser2025}{}}%
Visser, Coen, Alexander Heinlein, and Bianca Giovanardi. 2025.
{``PACMANN: Point Adaptive Collocation Method for Artificial Neural
Networks.''} arXiv preprint. \url{https://arxiv.org/abs/2411.19632v2}.

\leavevmode\vadjust pre{\hypertarget{ref-wang2021}{}}%
Wang, Sifan, Yujun Teng, and Paris Perdikaris. 2021. {``Understanding
and Mitigating Gradient Flow Pathologies in Physics-Informed Neural
Networks.''} \emph{SIAM Journal on Scientific Computing} 43 (January):
A3055--81. \url{https://doi.org/10.1137/20m1318043}.

\leavevmode\vadjust pre{\hypertarget{ref-wu2023}{}}%
Wu, Chenxi, Min Zhu, Qinyang Tan, Yadhu Kartha, and Lu Lu. 2023. {``A
Comprehensive Study of Non-Adaptive and Residual-Based Adaptive Sampling
for Physics-Informed Neural Networks.''} \emph{Computer Methods in
Applied Mechanics and Engineering} 403 (January): 115671.
\url{https://doi.org/10.1016/j.cma.2022.115671}.

\end{CSLReferences}

\end{document}
